\documentclass[10pt,a4paper]{amsart}
\usepackage[margin=19mm]{geometry}
\usepackage{amsmath,amssymb,mathtools}
\usepackage[scr=rsfs,cal=euler]{mathalfa}
\usepackage{xfrac}
\usepackage[unicode,hidelinks,bookmarks=false]{hyperref}
\newcommand{\ie}{i.e.\ }
\newcommand{\cf}{cf.\ }
\newcommand{\resp}{respectively\ }
\newcommand{\Subsection}{\subsection}
\providecommand{\nobreakdash}{\nobreak\hbox{-}}
\setlength{\emergencystretch}{2em}
\multlinegap0pt
\usepackage{bm}
\usepackage{bbm}
\usepackage{dsfont}
\usepackage{xspace}
\usepackage{cancel}
\usepackage{mathtools}
\usepackage{amsthm}
\makeatletter
\@ifundefined{definition}{\newtheorem{definition}{Definition}}{}
\@ifundefined{remark}{\newtheorem{remark}[definition]{Remark}}{}
\@ifundefined{theorem}{\newtheorem{theorem}[definition]{Theorem}}{}
\makeatother
\title[Common Fixed Points for Weakly Compatible Mappings via Bivar]{Common Fixed Points for Weakly Compatible Mappings via Bivariate Auxiliary Functions}
\author{Babu G.V.R., Alemayehu Negash, Meaza Bogale}
\date{Source: arXiv:2508.08260v1 (CC BY 4.0)}
\begin{document}
\maketitle

\noindent\emph{Source and licence.} Common Fixed Points for Weakly Compatible Mappings via Bivariate Auxiliary Functions. Version arXiv:2508.08260v1 (CC BY 4.0), distributed under the Creative Commons Attribution 4.0 International licence, as recorded in the arXiv licence metadata for this exact version and shown on the abstract page https://arxiv.org/abs/2508.08260v1. Full rights notice: licenses/arxiv-2508.08260-CC-BY-4.0.txt.\par\smallskip

\noindent\textit{Editorial scope.} Counted unit: the Theorem whose source label is \texttt{thm:main} in the article (source0.tex lines 202--214, source label \texttt{thm:main}), delivered together with the article's own complete original proof of it (source0.tex lines 216--463, one \texttt{proof} environment of 248 lines). No mathematics is written, repaired or re-derived: statement and proof are the article's own text line for line. Required context retained uncounted: rem:1.3 (lines 93--95). Every definition, display and label of those lines is retained unchanged. Omitted from this package and not redistributed: 1--92, 96--201, 215--215, 464--689. Nothing inside a retained excerpt is omitted or altered: every retained body is delivered exactly as the article's lines are written, so no omission receipt of any kind accompanies this package. Citations: the retained excerpts carry no citation command at all, so no bibliography entry of the article is transcribed and no reference is redistributed. Reference adaptation: none was needed and none was made; every reference inside the delivered unit resolves inside it, so no unresolved reference or citation is produced. Printed slips kept literal: no printed word, symbol, hypothesis or number of the counted statement or of its proof was altered. Layout and class: the article's own class is \texttt{article} with packages including amsmath, amssymb, graphicx, amsthm, mathrsfs, mathptmx, mathastext, mathtools, inputenc, enumitem, hyperref, geometry; the wrapper is the lane's pinned \texttt{amsart} editorial wrapper at 10pt on A4, which restates the theorem-like environments the retained text needs and adds the article's own macro definitions that the retained text needs (none), with extra wrapper packages (bm, bbm, dsfont, xspace, cancel, mathtools). The delivered numbering is the unit's own. Assets: no retained span contains a figure environment, a graphics-inclusion command, a PDF-page inclusion, a table or a TikZ picture; no image file is copied into this package, the package declares no asset dependency and no image file is redistributed with it. Licence: version arXiv:2508.08260v1 of this article is distributed under the Creative Commons Attribution 4.0 International licence, as recorded in the arXiv licence metadata for this exact version and shown on the abstract page https://arxiv.org/abs/2508.08260v1. A full rights notice is delivered at licenses/arxiv-2508.08260-CC-BY-4.0.txt. Characters: every retained excerpt is pure ASCII, so the delivered engine, XeLaTeX with its default Latin Modern text font, reports no missing character.
\begin{remark} \label{rem:1.3}
If $\psi(s,t)=0$, then $s=0$.
\end{remark}

\begin{theorem}\label{thm:main}
Let $(X,d)$ be a metric space and $K$ be a nonempty subset of $X$. Let $A,B,S,T$ be selfmaps on $K$. Suppose there exist $\psi\in\Psi$, $\phi\in\Phi$ and such that for all $x,y\in K$:

\begin{align} \label{eq:2.1}
    \psi(d(Sx,Ty),d(Sx,Ty)) \leq \phi(&\max\{\psi(d(Ax,By),d(Ax,Sx)),\psi(d(Ax,By),d(By,Ty)),\nonumber \\
 &\psi(d(Ax,Sx),d(By,Ty)), \psi(d(By,Ty),d(Ax,Sx)),\\
 &\min\{\psi(d(By,Sx),d(Ax,Sx)),
\psi(d(Ax,Ty),d(By,Ty))\},  \nonumber\\
&\min\{\psi(d(By,Sx),d(By,Ty)), 
\psi(d(Ax,Ty),d(Ax,Sx))\}\}). \nonumber
\end{align}
Assume also that the pairs of selfmaps $(A,S)$ and $(B,T)$ are weakly compatible; $\overline{T(K)}\subset A(K)$ and $\overline{S(K)}\subset B(K)$. If either $\overline{T(K)}$ or $\overline{S(K)}$ or $A(K)$ or $B(K)$ is complete, then $A$, $B$, $S$ and $T$ have a unique common fixed point.
\end{theorem}

\begin{proof}
Let $x_{0}\in K$. Since $T(K)\subseteq\overline{T(K)}\subseteq A(K)$, we can find $x_{1}\in K$ such that $Tx_{0}=Ax_{1}$ and since $S(K)\subseteq\overline{S(K)}\subseteq B(K)$, we can find $x_{2}\in K$ such that $Sx_{1}=Bx_{2}$. Continuing this way,  We construct sequences  $\{x_{n}\}$ and $\{y_{n}\}$ in $K$ inductively as follows:
\[
y_{2n} = Tx_{2n} = Ax_{2n+1} \quad \text{and} \quad y_{2n+1} = Sx_{2n+1} = Bx_{2n+2}
\]
for $n=0,1,2,\cdots$.
Now we claim that the sequence $\{y_{n}\}_{n=1}^{\infty}$ is Cauchy in $K$.

\textbf{Case (i):} $y_{n}=y_{n+1}$ for some $n$.

Without any loss of generality, we assume that $n$ is even. Then $n=2m$ for some $m\in \mathbb{Z}_{+}$. So $y_{2m}=y_{2m+1}$.

Thus, we have
\begin{align*}
\psi(d(y_{n+1},y_{n+2}),d(y_{n+1},y_{n+2})) &= \psi(d(y_{2m+1},y_{2m+2}),d(y_{2m+1},y_{2m+2})) \\
&= \psi(d(Sx_{2m+1},Tx_{2m+2}),d(Sx_{2m+1},Tx_{2m+2})) \\
&\leq \phi(\max\{\psi(d(Ax_{2m+1},Bx_{2m+2}),d(Ax_{2m+1},Sx_{2m+1})), \\
&\quad \psi(d(Ax_{2m+1},Bx_{2m+2}),\psi(d(Bx_{2m+2}, Tx_{2m+2})),\\
& \psi(d(Ax_{2m+1}, Sx_{2m+1}), d(Bx_{2m+2}, Tx_{2m+2})), \\
&\psi(d(Bx_{2m+2}, Tx_{2m+2}), d(Ax_{2m+1}, Sx_{2m+1})), \\
&\min\{\psi(d(Bx_{2m+2}, Sx_{2m+1}), d(Ax_{2m+1}, Sx_{2m+1})), \\
&\quad \psi(d(Ax_{2m+1}, Tx_{2m+2}), d(Bx_{2m+2}, Tx_{2m+2}))\}, \\
&\min\{\psi(d(Bx_{2m+2}, Sx_{2m+1}), d(Bx_{2m+2}, Tx_{2m+2})), \\
&\quad \psi(d(Ax_{2m+1}, Tx_{2m+2}), d(Ax_{2m+1}, Sx_{2m+1}))\} \\
&= \phi(\max\{\psi(d(y_n, y_{n+1}), d(y_n, y_{n+1})), \\
&\quad \psi(d(y_n, y_{n+1}), d(y_{n+1}, y_{n+2})), \quad \psi(d(y_n, y_{n+1}), d(y_{n+1}, y_{n+2})),\\
&\quad \psi(d(y_{n+1}, y_{n+2}), d(y_n, y_{n+1})), \\
&\quad \min\{\psi(d(y_{n+1}, y_{n+1}), d(y_n, y_{n+1})), \quad \psi(d(y_n, y_{n+2}), d(y_{n+1}, y_{n+2}))\},\\
&\quad \min\{\psi(d(y_{n+1}, y_{n+1}), d(y_{n+1}, y_{n+2})), \quad \psi(d(y_n, y_{n+2}), d(y_n, y_{n+1}))\}\\
&\leq \phi(\max\{0, \psi(0, d(y_{n+1}, y_{n+2})), \psi(d(y_{n+1}, y_{n+2}), 0)\}) \\
&\leq \phi(\psi(d(y_{n+1}, y_{n+2}), d(y_{n+1}, y_{n+2}))).
\end{align*}

This implies that $\psi(d(y_{n+1}, y_{n+2}), d(y_{n+1}, y_{n+2})) = 0$ and hence $d(y_{n+1}, y_{n+2}) = 0$ (By Remark \ref{rem:1.3}).

Hence, $y_{n+1} = y_{n+2}$. Repeating this procedure inductively, we show that $y_n = y_{n+k}$ for $k \geq 1$. Hence, $\{y_m\}_{m\geq n}$ is a constant sequence and therefore is Cauchy in $K$.

\textbf{Case (ii):} $y_n \neq y_{n+1}$, $n = 1, 2, \dotsc$.

Let $\alpha_n = d(y_n, y_{n+1})$ for all $n = 0, 1, 2, \dotsc$.

Now consider
\begin{align*}
\psi(\alpha_{2n+1}, \alpha_{2n+1}) &= \psi(d(y_{2n+1}, y_{2n+2}), d(y_{2n+1}, y_{2n+2})) \\
&= \psi(d(Sx_{2n+1}, Tx_{2n+2}), d(Sx_{2n+1}, Tx_{2n+2})) \\
&\leq \phi(\max\{\psi(d(Ax_{2n+1}, Bx_{2n+2}), d(Ax_{2n+1}, Sx_{2n+1})), \psi(d(Ax_{2n+1}, Bx_{2n+2}), d(Bx_{2n+2}, Tx_{2n+2})),\\
&\quad \quad \psi(d(Ax_{2n+1}, Sx_{2n+1}), d(Bx_{2n+2}, Tx_{2n+2})), \psi(d(Bx_{2n+2}, Tx_{2n+2}), d(Ax_{2n+1}, Sx_{2n+1})),\\
&\quad \quad \min\{\psi(d(Bx_{2n+2}, Sx_{2n+1}), d(Ax_{2n+1}, Sx_{2n+1})), \psi(d(Ax_{2n+1}, Tx_{2n+2}), d(Bx_{2n+2}, Tx_{2n+2}))\}, \\
&\quad \quad \min\{\psi(d(Bx_{2n+2}, Sx_{2n+1}), d(Bx_{2n+2}, Tx_{2n+2})), \psi(d(Ax_{2n+1}, Tx_{2n+2}), d(Ax_{2n+1}, Sx_{2n+1}))\}\\
&= \phi(\max\{\psi(d(y_{2n}, y_{2n+1}), d(y_{2n}, y_{2n+1})), \psi(d(y_{2n}, y_{2n+1}), d(y_{2n+1}, y_{2n+2})),\\
&\quad \quad \psi(d(y_{2n}, y_{2n+1}), d(y_{2n+1}, y_{2n+2})), \psi(d(y_{2n+1}, y_{2n+2}), d(y_{2n}, y_{2n+1})), \\
&\quad \quad \min\{\psi(d(y_{2n+1}, y_{2n+1}), d(y_{2n}, y_{2n+1})), \psi(d(y_{2n}, y_{2n+2}), d(y_{2n+1}, y_{2n+2}))\}, \\
&\quad \quad \min\{\psi(d(y_{2n+1}, y_{2n+1}), d(y_{2n+1}, y_{2n+2})), \psi(d(y_{2n}, y_{2n+2}), d(y_{2n}, y_{2n+1}))\} \\
&\leq \phi(\max\{\psi(\alpha_{2n}, \alpha_{2n}), \psi(\alpha_{2n}, \alpha_{2n+1}), \psi(\alpha_{2n+1}, \alpha_{2n}),\psi(0, \alpha_{2n}), \psi(0, \alpha_{2n+1})\}).
\end{align*}

Now if $\alpha_{2n} < \alpha_{2n+1}$, then $\psi(\alpha_{2n+1}, \alpha_{2n+1}) \leq \phi(\psi(\alpha_{2n+1}, \alpha_{2n+1}))$, which is a contradiction. Hence, $\alpha_{2n+1} \leq \alpha_{2n}$ and
\begin{equation}\label{eq:2.1.2}
\psi(\alpha_{2n+1}, \alpha_{2n+1}) \leq \phi(\psi(\alpha_{2n}, \alpha_{2n})), \quad n = 0, 1, 2, \dotsc.
\end{equation}

Similarly, we can show that
\begin{equation}\label{eq:2.1.3}
\psi(\alpha_{2n}, \alpha_{2n}) \leq \phi(\psi(\alpha_{2n-1}, \alpha_{2n-1})), \quad n = 1, 2, 3, \dotsc.
\end{equation}

Hence, from \eqref{eq:2.1.2} and \eqref{eq:2.1.3}, we get
\begin{equation}\label{eq:2.1.4}
\psi(\alpha_{n+1}, \alpha_{n+1}) \leq \phi(\psi(\alpha_{n}, \alpha_{n})), \quad n = 0, 1, 2, \dotsc.
\end{equation}

Since $\psi$ is monotone increasing in each of its variables, we have
\begin{equation}\label{eq:2.1.6}
\alpha_{2n+1} \leq \alpha_{2n}, \quad n = 0, 1, 2, \dotsc.
\end{equation}

Thus, the sequence $\{\alpha_n\}$ decreases, say, to $\alpha (\geq 0)$. Since $\phi$ and $\psi$ are continuous, letting $n \to \infty$, from \eqref{eq:2.1.4}, we get
\[
\psi(\alpha, \alpha) \leq \phi(\psi(\alpha, \alpha)),
\]
so that $\psi(\alpha, \alpha) = 0$ and hence $\alpha = 0$. Therefore,
\begin{equation}\label{eq:2.1.7}
\lim_{n\to\infty} d(y_n, y_{n+1}) = 0.
\end{equation}

We now claim that $\{y_n\}$ is Cauchy. By \eqref{eq:2.1.6} and \eqref{eq:2.1.7}, it is sufficient to show that $\{y_{2n}\}$ is Cauchy. Otherwise, there exists an $\epsilon > 0$ and there exist sequences $\{m_k\}$ and $\{n_k\}$ with $m_k > n_k > k$ such that
\begin{equation}\label{eq:2.1.8}
d(y_{2m_k}, y_{2n_k}) \geq \epsilon \quad \text{and} \quad d(y_{2m_k-2}, y_{2n_k}) < \epsilon.
\end{equation}

Hence,
\begin{equation}\label{eq:2.1.9}
\epsilon \leq \liminf_{k\to\infty} d(y_{2m_k}, y_{2n_k}).
\end{equation}

Now for each positive integer $k$, by triangle inequality we get,
\begin{equation*}
d(y_{2m_k}, y_{2n_k}) \leq d(y_{2m_k}, y_{2m_k-1}) + d(y_{2m_k-1}, y_{2m_k-2}) + d(y_{2m_k-2}, y_{2n_k}).
\end{equation*}

On taking limit supremum of both sides, as $k \to \infty$, we get
\begin{equation}\label{eq:2.1.10}
\limsup_{k\to\infty} d(y_{2m_k}, y_{2n_k}) \leq \epsilon.
\end{equation}

Hence from \eqref{eq:2.1.9} and \eqref{eq:2.1.10}, we have
\begin{equation}\label{eq:2.1.11}
\lim_{k\to\infty} d(y_{2m_k}, y_{2n_k}) = \epsilon.
\end{equation}

Now
\begin{align*}
d(y_{2m_k}, y_{2n_k-1}) &\leq d(y_{2m_k}, y_{2n_k}) + d(y_{2n_k}, y_{2n_k-1});
\end{align*}
On taking limit supremum, as $k \to \infty$, we get
\begin{equation}\label{eq:2.1.12}
\limsup_{k\to\infty} d(y_{2m_k}, y_{2n_k-1}) \leq \epsilon.
\end{equation}

Again we have
\begin{align*}
d(y_{2m_k}, y_{2n_k}) &\leq d(y_{2m_k}, y_{2n_k-1}) + d(y_{2n_k-1}, y_{2n_k});
\end{align*}
On taking limit infimum, as $k \to \infty$, we get
\begin{equation}\label{eq:2.1.13}
\epsilon \leq \liminf_{k\to\infty} d(y_{2m_k}, y_{2n_k-1}).
\end{equation}

From \eqref{eq:2.1.12} and \eqref{eq:2.1.13}, we have
\begin{equation}\label{eq:2.1.14}
\lim_{k\to\infty} d(y_{2m_k}, y_{2n_k-1}) = \epsilon.
\end{equation}

Similarly, we can show that
\begin{equation}\label{eq:2.1.15}
\lim_{k\to\infty} d(y_{2m_k+1}, y_{2n_k}) = \epsilon
\end{equation}
and
\begin{equation}\label{eq:2.1.16}
\lim_{k\to\infty} d(y_{2m_k+1}, y_{2n_k-1}) = \epsilon.
\end{equation}

Now consider
\begin{align*}
\psi(d(y_{2m_k+1}, y_{2n_k}), d(y_{2m_k+1}, y_{2n_k})) &= \psi(d(Sx_{2m_k+1}, Tx_{2n_k}), d(Sx_{2m_k+1}, Tx_{2n_k})) \\
&\leq \phi(\max\{\psi(d(Ax_{2m_k+1}, Bx_{2n_k}), d(Ax_{2m_k+1}, Sx_{2m_k+1})), \\
&\quad \psi(d(Ax_{2m_k+1}, Bx_{2n_k}), d(Bx_{2n_k}, Tx_{2n_k})), \\
&\quad \psi(d(Ax_{2m_k+1}, Sx_{2m_k+1}), d(Bx_{2n_k}, Tx_{2n_k})), \\
&\quad \psi(d(Bx_{2n_k}, Tx_{2n_k}), d(Ax_{2m_k+1}, Sx_{2m_k+1})), \\
&\quad \min\{\psi(d(Bx_{2n_k}, Sx_{2m_k+1}), d(Ax_{2m_k+1}, Sx_{2m_k+1})), \\
&\quad \quad \psi(d(Ax_{2m_k+1}, Tx_{2n_k}), d(Bx_{2n_k}, Tx_{2n_k}))\}, \\
&\quad \min\{\psi(d(Bx_{2n_k}, Sx_{2m_k+1}), d(Bx_{2n_k}, Tx_{2n_k})), \\
&\quad \quad \psi(d(Ax_{2m_k+1}, Tx_{2n_k}), d(Ax_{2m_k+1}, Sx_{2m_k+1}))\} \\
&= \phi(\max\{\psi(d(y_{2m_k}, y_{2n_k-1}), d(y_{2m_k}, y_{2m_k+1})), \psi(d(y_{2m_k}, y_{2n_k-1}), d(y_{2n_k-1}, y_{2n_k})),\\
&\quad \min\{\psi(d(y_{2n_k-1},y_{2m_k+1}),d(y_{2m_k},y_{2m_k+1})), \psi(d(y_{2m_k},y_{2n_k}),d(y_{2n_k-1},y_{2n_k}))\},\\
&\quad \min\{\psi(d(y_{2n_k-1},y_{2m_k+1}),d(y_{2n_k-1},y_{2n_k})),\psi(d(y_{2m_k},y_{2n_k}),d(y_{2m_k},y_{2m_k+1}))\}
\end{align*}

Letting $n\to\infty$, by the continuity of $\phi$ and $\psi$, and from \eqref{eq:2.1.7}, \eqref{eq:2.1.11}, \eqref{eq:2.1.14}, \eqref{eq:2.1.15} and \eqref{eq:2.1.16}, we get

\begin{align*}
\psi(\varepsilon,\varepsilon) &\leq \phi(\max\{\psi(\varepsilon,0),\psi(\varepsilon,0),\psi(0,0), \psi(0,0), \\
&\quad \min\{\psi(\varepsilon,0),\psi(\varepsilon,0)\}, \min\{\psi(\varepsilon,0),\psi(\varepsilon,0)\}\}) \\
&= \phi(\psi(\varepsilon,0)) \\
&\leq \phi(\psi(\varepsilon,\varepsilon)), \text{ a contradiction}.
\end{align*}

Hence, $\{y_{2n}\}$ is a Cauchy sequence in $K$ and hence $\{y_n\}$ is Cauchy in $K$.

Hence, the subsequences $\{Tx_{2n}\}$, $\{Sx_{2n+1}\}$, $\{Ax_{2n}\}$ and $\{Bx_{2n+2}\}$ are Cauchy in $K$.

Assume that $\overline{T(K)}$ is complete. Then there exists a $z\in\overline{T(K)}$ such that
\begin{equation}\label{eq:2.1.17}
\lim_{n\to\infty} Tx_{2n} = z
\end{equation}
and hence,
\begin{equation}\label{eq:2.1.18}
\lim_{n\to\infty} Ax_{2n+1} = \lim_{n\to\infty} Bx_{2n+2} = \lim_{n\to\infty} Sx_{2n+1} = z
\end{equation}
and $z\in\overline{S(K)}$.

Since $\overline{T(K)} \subset A(K)$ and $\overline{S(K)} \subset B(K)$, there exist $u,v\in K$ such that
\begin{equation}\label{eq:2.1.19}
z = Au \text{ and } z = Bv.
\end{equation}

Now taking $x = u$ and $y = x_{2n}$ in \eqref{eq:2.1}, we have
\begin{align*}
\psi(d(Su,Tx_{2n}),d(Su,Tx_{2n})) &\leq \phi(\max\{\psi(d(Au,Bx_{2n}),d(Au,Su)),\psi(d(Au,Bx_{2n}), d(Bx_{2n},Tx_{2n})),  \\
&\quad \psi(d(Au,Su),d(Bx_{2n},Tx_{2n})), \psi(d(Bx_{2n},Tx_{2n}),d(Au,Su)),\\
&\quad \min\{\psi(d(Bx_{2n},Su),d(Au,Su)), \psi(d(Au,Tx_{2n}),d(Bx_{2n},Tx_{2n}))\},\\
&\quad \min\{\psi(d(Bx_{2n},Su), d(Bx_{2n},Tx_{2n})), \psi(d(Au,Tx_{2n}),d(Au,Su))\})
\end{align*}

Letting $n\to\infty$, by the continuity and monotone increasing property of $\phi$ and $\psi$ and using \eqref{eq:2.1.17}, \eqref{eq:2.1.18} and \eqref{eq:2.1.19}, we get
\[
\psi(d(Su,z),d(Su,z)) \leq \phi(\psi(d(Su,z),d(Su,z))).
\]
This implies that $\psi(d(Su,z),d(Su,z)) = 0$ and hence by Remark \ref{rem:1.3} we have $d(z,Su) = 0$.

Hence, $z = Su$.

Similarly in \eqref{eq:2.1}, by taking $x = x_{2n+1}$ and $y = v$ and letting $n\to\infty$, we get $z = Tv$.

Hence, $Su = Au = z = Tv = Bv$.

Since the pairs of mappings $(A,S)$ and $(B,T)$ are weakly compatible mappings, we have $ASu = SAu$ and $BTv = TBv$ and hence
\[
Az = Sz \text{ and } Bz = Tz.
\]

Hence, $z$ is the point of coincidence of $A$ and $S$, and $B$ and $T$.

We now claim that $Az = Sz = z$. Consider
\begin{align*}
\psi(d(Sz,z),d(Sz,z)) &= \psi(d(Sz,Tu),d(Sz,Tu)) \\
&\leq \phi(\max\{\psi(d(Az,Bu),d(Az,Sz)), \psi(d(Az,Bu),d(Bu,Tu)), \\
&\quad \psi(d(Az,Sz),d(Bu,Tu)), \psi(d(Bu,Tu),d(Az,Sz)), \\
&\quad \min\{\psi(d(Bu,Sz),d(Az,Sz)), \psi(d(Az,Tu),d(Bu,Tu))\}, \\
&\quad \min\{\psi(d(Bu,Sz),d(Bu,Tu)), \psi(d(Az,Tu),d(Az,Sz))\}\}) \\
&= \phi(\psi(d(Sz,z),0)) \\
&\leq \phi(\psi(d(Sz,z),d(Sz,z))).
\end{align*}

This implies that $\psi(d(Sz,z),d(Sz,z)) = 0$ and hence $d(Sz,z) = 0$. Therefore,
\[
Sz = z.
\]
Hence, $Az = Sz = z$.

Similarly, by taking $x = u$ and $y = z$ in \eqref{eq:2.1} we get $Tz = z$. Hence,
\[
Bz = Tz = z.
\]

Thus, $z$ is a common fixed point of $A$, $B$, $S$ and $T$. The uniqueness of $z$ follows from the inequality \eqref{eq:2.1}.

Similarly, by symmetry we prove for the case if $\overline{S(K)}$ is complete.

Now assume $A(K)$ is complete. As $\overline{T(K)} \subset A(K)$ and $\overline{T(K)}$ is closed, $\overline{T(K)}$ is complete. Hence by the above argument, there exists a point $z \in \overline{T(K)}$ such that
\[
Az = Bz = Sz = Tz = z.
\]

Similarly, if we assume $B(K)$ is complete, the proof follows by symmetry.

Hence the theorem.
\end{proof}

\end{document}
